% =====================================================================
%  CONFIGURACIÓN — paquetes, página, interlineado y estilos
%  Formato basado en la plantilla de tesis del CISE.
%  Un cambio aquí afecta a toda la tesis.
% =====================================================================

% ----------------------------- IDIOMA --------------------------------
\usepackage[utf8]{inputenc}
\usepackage[T1]{fontenc}
\usepackage[american, spanish, es-noquoting, es-nodecimaldot, es-tabla,
            es-lcroman]{babel}
% es-tabla: "Tabla" en lugar de "Cuadro"; es-nodecimaldot: 0.5 y no 0'5
% es-lcroman: páginas romanas en minúscula normal (ii, iii); sin esta
%             opción, babel las escribe en versalitas
\usepackage{csquotes}

% Los anexos se titulan "Anexo A" en lugar de "Apéndice A"
\addto\captionsspanish{\renewcommand{\appendixname}{Anexo}}

% ----------------------------- FUENTE --------------------------------
% Palatino (TeX Gyre Pagella) en el texto y en todas las ecuaciones,
% incluidos los símbolos (=, +, \sum, \int, \mathcal...). newpxmath carga
% amsmath y los símbolos de amssymb, así que no cargue amssymb aparte.
% amsmath y amsthm deben cargarse ANTES que newpxmath.
\usepackage{amsmath, amsthm}
\usepackage{newpxtext}
\usepackage{newpxmath}
\usepackage{microtype}

% ------------------------ GEOMETRÍA DE PÁGINA ------------------------
\usepackage[paper=letterpaper, centering, includehead]{geometry}
\geometry{tmargin=2cm, bmargin=2.5cm, lmargin=2cm, rmargin=2cm}
\geometry{head=0cm, foot=0cm, headsep=0cm, footskip=0.75cm}
\geometry{columnsep=0.5cm}

% ---------------------- PÁRRAFOS E INTERLINEADO ----------------------
\setlength{\parskip}{0.3cm}         % Espacio entre párrafos
\setlength{\parindent}{0cm}         % Sin sangría
\usepackage{setspace}
\setstretch{1.3}                    % Interlineado (1.5 = \onehalfspacing)

% ------------------------ ESTILO DE CAPÍTULOS ------------------------
\usepackage[Glenn]{fncychap}
\ChNameVar{\bfseries\Large\rmfamily}
\ChNumVar{\Huge}
\ChTitleVar{\bfseries\large\rmfamily}
\ChRuleWidth{1pt}

% ---------------------- ENCABEZADOS Y PIES ---------------------------
\usepackage{fancyhdr}
\pagestyle{fancy}
\fancyhf{}                          % Limpia encabezado y pie
\fancyfoot[R]{\thepage}             % Número de página abajo a la derecha
\renewcommand{\headrulewidth}{0pt}
\renewcommand{\footrulewidth}{0.5pt}
% Primera página de cada capítulo: mismo pie
\fancypagestyle{plain}{%
  \fancyhf{}\fancyfoot[R]{\thepage}%
  \renewcommand{\headrulewidth}{0pt}\renewcommand{\footrulewidth}{0.5pt}}

% --------------------------- MATEMÁTICAS -----------------------------
\newtheorem{teorema}{Teorema}[chapter]
\newtheorem{corolario}[teorema]{Corolario}
\newtheorem{lema}[teorema]{Lema}
\newtheorem{proposicion}[teorema]{Proposición}
\theoremstyle{definition}
\newtheorem{definicion}{Definición}[chapter]
\newtheorem{ejemplo}{Ejemplo}[chapter]

% ------------------------ TABLAS Y FIGURAS ---------------------------
\usepackage{graphicx}
\graphicspath{{figuras/}}           % Las imágenes se buscan en figuras/
\usepackage{booktabs}
\usepackage{multirow}
\usepackage{longtable}
\usepackage{threeparttable}
\usepackage{caption}
\captionsetup{font=small, labelfont=bf}
\captionsetup[table]{skip=10pt}
\usepackage{subcaption}
\usepackage{float}
\usepackage[dvipsnames]{xcolor}

% ------------------------------ LISTAS -------------------------------
\usepackage{enumitem}
\setlist{itemsep=0.1cm, topsep=0.1cm}

% --------------------------- OTROS -----------------------------------
\usepackage{multicol}
\usepackage[nohyperlinks]{acronym}  % Siglas: \ac{INEGI} (ver codigo/siglas.tex)
\usepackage{lipsum}                 % Texto de relleno: borrar al terminar
\SetLipsumLanguage{spanish}         % Evita el aviso por el idioma "latin"

% --------------------------- MARCA DE AGUA ---------------------------
% Para versiones de revisión, descomente las siguientes líneas.
%\usepackage{draftwatermark}
%\SetWatermarkText{BORRADOR}
%\SetWatermarkLightness{0.93}
%\SetWatermarkScale{0.75}

% -------------------------- BIBLIOGRAFÍA -----------------------------
% APA 7 (paquete biblatex-apa). No agregue opciones como maxcitenames:
% style=apa ya aplica las reglas de APA (et al., número de autores, etc.).
% Para un autor-año genérico en lugar de APA, cambie la línea por:
%   \usepackage[backend=biber, style=authoryear, maxcitenames=2,
%               maxbibnames=10, uniquename=false, dashed=false]{biblatex}
\usepackage[backend=biber, style=apa]{biblatex}
\addbibresource{referencias.bib}

% ------------------------- COMANDOS PROPIOS --------------------------
\newcommand{\E}{\mathbb{E}}         % Esperanza
\newcommand{\R}{\mathbb{R}}         % Números reales
\newcommand{\ingles}[1]{\foreignlanguage{american}{\emph{#1}}} % Términos en inglés

% ---------------- ENLACES Y REFERENCIAS CRUZADAS ---------------------
% Siempre al final del preámbulo y en este orden.
\usepackage[hidelinks, bookmarksnumbered]{hyperref}
\usepackage[spanish]{cleveref}
\crefname{chapter}{capítulo}{capítulos}
\crefname{section}{sección}{secciones}
\crefname{equation}{ecuación}{ecuaciones}
\crefname{table}{tabla}{tablas}
\crefname{figure}{figura}{figuras}
\crefname{appendix}{anexo}{anexos}
